\section{Game-Based Protocol Security}
\label{sec:security}
The statistical proofs do not establish message authenticity, and the
historical Lean model supplies authentication as a Boolean premise. This
section gives an explicit computational experiment connecting signature and
hash security to a narrower protocol property: correspondence between an
accepted envelope and previously issued content. It also proves at-most-once
effects under atomic state management and states the additional premises
needed to combine authenticity with privacy assurance. These are written
proofs of the specified construction, not new kernel-checked results or
proofs of the Python implementation. The signature reduction is a standard
application of chosen-message unforgeability, not a new cryptographic primitive
\citep{goldwasser1988signatures}.

\paragraph{Security framework}
This section uses a game-based authenticity definition with explicit
adversary oracles and a bad-event winning condition \citep{bellare2006games}.
S1--S3 are not universal-composability (UC) results. UC requires an ideal
functionality, a simulator and a real/ideal argument against the allowed
environments \citep{canetti2001uc,canetti2025uc}; an idealized registry or
a union of failure events alone is insufficient.
Section~\ref{sec:uc} separately supplies those objects for a restricted
public-record service in a static authenticated-channel model.
The scientific premises remain external requirements, and concrete
transport substitution is conditional. S3 is a joint failure-event
bound within its own experiment, not a UC composition theorem.

\subsection{Game and adversary model}
\label{sec:security-game}
Fix a security parameter $\lambda$, a polynomial execution horizon, and $N$
protected, enrolled public keys independently generated by a signature
scheme $\Pi=(\mathrm{KeyGen},\mathrm{Sign},\mathrm{Verify})$. Those private
keys remain uncompromised throughout the experiment. Enrollment binds the
full public key to issuer and permitted roles; short display identifiers
cannot select an attacker-supplied replacement key. Other identities may be
adversarial, but their signatures cannot satisfy a protected issuer's
requirement. Policy, trust-directory and clock queries use the authoritative
values of the state model, not network-supplied assertions.

The probabilistic polynomial-time adversary controls message delivery,
ordering, delay, replay and untrusted object storage. It may adaptively
request signatures for any syntactically supported typed messages under
the protected keys. The signing oracle need not judge whether their
scientific claims are true. The adversary sees all issued messages and
signatures and may submit polynomially many verification attempts. It
cannot rewrite the verifier, authoritative state or protected secret keys
in this experiment. Those attacks require separate implementation and
infrastructure arguments.

Let $C$ be an injective encoding of the admitted typed values, including
explicit type tags and unambiguous lengths. Distinct field sequences cannot
share an encoding; malformed, ambiguous and unsupported inputs are rejected.
Canonical message equality means equality of these typed values, not equality
of irrelevant transport whitespace. The proof-level envelope $x$ binds
protocol domain/version/type, full issuer key and role/action, registry and
instance identity, message identifier, nonce, expected head, key/trust epoch,
issue/expiry times, predecessor identifiers/digests and complete payload.
It also binds an ordered inventory of named objects. Each captured object
byte string $b_\ell$ has typed commitment
\[
 d_\ell=H(C(\mathrm{object},\mathrm{kind}_\ell,
                 \mathrm{name}_\ell,b_\ell)).
\]
The inventory includes the full context, policy, model/interface manifests
and every transitive predecessor or evidence object used by the decision;
fields represented inline are part of $x$ itself. A finite acyclic inventory
is required. Object kind/name and inventory position prevent untyped
interchange of equal-looking references.

The issuance oracle captures and retains the original typed envelope and
all referenced bytes, computes their commitments, and returns
\[
 m=H(C(\mathrm{envelope},x)),\qquad
 \sigma\leftarrow\mathrm{Sign}(sk_i,m).
\]
It records $(i,x,(b_\ell)_\ell,m)$ in its issuance log. An arbitrary
attacker-supplied digest without a captured original preimage is not an
issuance record for the object-binding claim. Verification captures each
supplied object once, hashes and parses those same bytes, recomputes the
complete inventory and checks $\mathrm{Verify}(pk_i,m,\sigma)$ together
with the current A1 guards. No path is reopened after hashing.

Define $\mathrm{BadBind}$ to occur when a successful acceptance attributed
to a protected issuer has no earlier issuance-log entry with the same
typed envelope and every referenced byte string. This is an existence
correspondence, not a unique signing-event claim: identical signing requests
may occur more than once. Acceptance of a previously issued message with a
new signature is not itself a win.

\begin{figure*}[t]
\centering
\begin{minipage}{0.97\textwidth}
\hrule\smallskip
\normalsize
\textbf{Game $G_{\rm bind}^{\mathcal A}(1^\lambda)$: challenger and oracles.}
All query counts and total input length are polynomially bounded.
\begin{enumerate}
\item \textbf{Initialize.} Generate the $N$ protected key pairs; install their
full public keys and roles in the authoritative directory. Initialize
the specified state $\Sigma$ and an issuance log $L$. Any protected envelopes
preloaded into the initial state must also be generated through Issue and
recorded in $L$. Give $\mathcal A$ the public parameters, keys and oracle
interfaces. All protected signing, including outputs generated internally
by protocol/environment actions, invokes Issue, records its entry before
exposure and counts toward the signing-query budget. There is no corruption
oracle for protected keys.
\item \textbf{Issue$(i,u,B)$.} Reject unsupported typed fields or invalid/cyclic
object inventories. Capture $B=(b_\ell)_\ell$, construct envelope $x$ from
$u$ and the computed typed commitments, and set
$m=H(C(\mathrm{envelope},x))$.
Generate $\sigma\leftarrow\mathrm{Sign}(sk_i,m)$; append $(i,x,B,m)$ to $L$
before returning $(x,\sigma)$. This oracle does not check scientific truth,
consume a replay token, commit a predecessor or authorize a protocol step.
\item \textbf{TryAccept$(x',B',\sigma')$.} Capture the submitted inputs and
run the fixed A1 verifier against current $\Sigma$. Reject failed parsing,
signature or current-guard checks. For each envelope/object bundle actually
accepted under a protected key $i$, immediately test whether an earlier
entry of $L$ matches $i$, the typed envelope and every consumed reference
byte string. If none matches, stop and output $1$. Otherwise return the
verifier's ordinary result. A historical-receipt response validates no
unused replacement input and grants no new service; correspondence concerns
the content actually validated or recovered from the recorded history.
\item \textbf{Schedule$(a)$.} The adversary can select delivery order, delay
and the next requested protocol/environment action. Execute only a transition
admitted by the fixed state model, checking its authority and prerequisites;
otherwise refuse. The adversary cannot assign arbitrary authoritative state.
Issuance and verifier success do not bypass the transition rules. S2
separately requires atomic persistence for state-changing effects.
\item \textbf{Finish.} If the adversary halts, or its declared resource bound
is exhausted without a winning acceptance, output $0$. Later issuance cannot
repair an earlier win.
\end{enumerate}
\smallskip\hrule
\end{minipage}
\caption{The S1 authenticity game. Every adversarial envelope submitted through
a scheduled protocol action invokes the same TryAccept checks; Schedule is
not an alternative unchecked admission path.}
\label{fig:binding-game}
\end{figure*}

Figure~\ref{fig:binding-game} makes the existing experiment explicit.
Its advantage is the failure probability
\[
 \operatorname{Adv}^{\rm bind}_{\Pi,H}(\mathcal A;\lambda)
 =\Pr[G_{\rm bind}^{\mathcal A}(1^\lambda)=1].
\]
Probability includes key generation, signing, adversary and any modeled
environment randomness. This is not a bit-distinguishing game, so no
$1/2$ guessing baseline is subtracted. S1 quantifies over every admitted
polynomial-time adversary, including adaptive signing and verification
queries and polynomially many instances within this fixed interface.
That scope does not imply arbitrary-protocol UC composition.

\subsection{Security statements}
\begin{securitytheorem}[Envelope and object correspondence]
\label{sec-thm:binding}
For the experiment above there are signature-forgery and hash-collision
reductions such that
\[
 \operatorname{Adv}^{\rm bind}_{\Pi,H}(\mathcal A;\lambda)
 \le N\epsilon_{\rm sig}(\lambda)+\epsilon_H(\lambda).
\]
Here $\epsilon_{\rm sig}$ bounds single-key existential unforgeability under
adaptive chosen-message attack (EUF-CMA), and $\epsilon_H$ bounds collision
finding for the typed uses of $H$, at the resources of those reductions.
Their time is adversary time plus polynomial simulation, verification and
inventory-comparison cost; the signature reduction makes no more signing
queries than the original experiment. The collision reduction covers all
envelope and object domains together. No independence or per-verification
failure multiplication is assumed.
\end{securitytheorem}
The full proof appears in Appendix~\ref{app:security-proofs}. A different
envelope with the same signed digest, or a different referenced byte string
with the same commitment, supplies a hash collision. A new valid signed
digest supplies a signature forgery. This division is essential: signature
unforgeability alone does not prove that a hash reference names unique bytes.
The bounds are symbolic assumptions; this paper assigns no measured
failure probability to Ed25519, SHA-256 or a deployed key service.

\begin{securitytheorem}[At-most-once effect and current admission]
\label{sec-thm:replay}
Suppose each state-changing acceptance atomically checks and consumes a
signed replay token $(\mathrm{registry},\mathrm{issuer},\mathrm{epoch},
\mathrm{nonce})$, checks the current head and guards, and persists its
effect and receipt together. Each successful commit installs a fresh head
different from its predecessor. Consumed tokens are never removed and old
heads are never restored. Invalidated records cannot regain eligibility
without a newly checked applicable authorization. Registry, trust and
revocation updates and grants have one linearizable ordering.
Grant guards require a current ACTIVE
authorization, matching authorization/issuer epochs, exact context and
artifact/interface bindings, and strict live authorization/lease deadlines
at that linearization point. Then one token causes at most one new transition.
Two commits against one predecessor head cannot both succeed. A grant
linearized after revocation, suspension, epoch replacement or expiry fails
unless a fresh applicable authorization has been established.
\end{securitytheorem}
This is an additional sufficient concurrency requirement on the construction:
A1 specifies message-ID deduplication and A5 consumes commitment nonces;
S2 requires every state-changing caller to check and consume its signed
token atomically with the effect. Equivalently A1's signed message identifier,
with its registry/instance/issuer namespace, can be the persistent replay
token. A sequential A1 call alone does not establish this atomicity.
The result concerns new grants, not arrival of output from an earlier grant.
Returning a historical receipt produces no new transition or current service
authorization. Deduplicating signature bytes would not suffice: EUF-CMA
does not prohibit another valid signature on an already signed message.
Rollback resistance and faithful current-state reads are state-management
premises, not consequences of signature verification.

\begin{securitytheorem}[Joint authentication and ordinary-release failure bound]
\label{sec-thm:composition}
For grants, use the coupled execution that continues to the fixed horizon
after a correspondence failure, retaining an absorbing $\mathrm{BadBind}$
flag. Its flag probability equals the early-stopping game's win probability.
In this adaptive execution, assume the noncryptographic transition,
current-state and complete-mediation premises of
Theorem~\ref{thm:lifecycle} and Security Theorem~\ref{sec-thm:replay}.
Every issuer whose envelope is relied on to support a grant must belong
to the protected key set.
On executions without $\mathrm{BadBind}$, require faithful scientific
collection, applicable complete-channel evidence and the admission/reduction
premises of Theorem~\ref{thm:statistics}. Define its coverage events and
allocations on the entire experiment, including adversarial scheduling and
selection, with the stated conditional failure bounds. If every grant
requires ordinary RELEASE, then
\[
 \begin{split}
 &\Pr[\exists\text{ grant with a registered }
                       \theta_{ctk}>\tau_{ctk}]\\
 &\hspace{1em}\le \alpha+
              N\epsilon_{\rm sig}(\lambda)+\epsilon_H(\lambda).
 \end{split}
\]
\end{securitytheorem}
The proof uses event inclusion and a union bound, not independence or
conditioning the statistical guarantee on cryptographic success. Coverage
valid only in a benign, fixed-selection experiment cannot be inserted here.
The event excludes accepted optional residual risk and unregistered threats.
Authentic false measurements, undisclosed channels, bypassed output paths
and dishonest authorizers can violate the premises; this result assigns no
numerical bound to those failures.

\subsection{Attacks outside the cryptographic conclusion}
\label{sec:security-counterexamples}
\paragraph{Authentic false evidence}
An approved worker can sign fabricated counts claiming a low-risk channel.
The bytes and signature are genuine, so $\mathrm{BadBind}$ need not occur.
If downstream admission trusts the counts without justified collection or
other applicable evidence, an unsafe candidate may pass. Thus S1 proves
origin correspondence, not measurement truth or scientific adequacy.

\paragraph{Rollback and delayed use}
Restore a valid registry snapshot from before a nonce was consumed or a
release revoked. Every retained signature can verify, yet a verifier using
only that snapshot can admit the old action again. No signature forgery or
hash collision is necessary. S2 requires rollback-resistant authoritative
state; a fresh externally trusted checkpoint is one possible supporting
mechanism, not the only implementation. Likewise, a token checked before
expiry can be delivered after expiry; S2 proves a property at the grant
linearization point, not a physical delivery deadline.

\paragraph{Compromise and omitted observations}
After obtaining a protected signing key, an attacker can issue arbitrary
valid envelopes. The no-compromise premise then fails; a key-revocation
check helps only once current trusted state records the change.
Separately, two correctly signed releases can have harmless marginals and
reveal a secret jointly, as in the XOR construction in Appendix~\ref{app:gaps}.
Authentication supplies neither portfolio completeness nor confidentiality.
Hashing low-entropy confidential content is also not encryption.

\paragraph{Implementation and novelty boundary}
The construction specifies A1's full domain-separated envelope and typed
object commitments. Python's generic \code{sign\_canonical} instead signs
canonical JSON directly, and the existing signing helpers do not jointly
implement this entire enrollment, envelope, replay and current-state service.
Its shortened key identifier is a display/routing value, not a substitute
for pinning the full approved public key. The separate SQLite reference
trusts its gate/authority inputs and does not verify these signatures.
Existing Lean declarations use a symbolic authentication flag; neither
they nor Python regression tests prove S1 or S3. Ed25519's design and
specification are established external work
\citep{bernstein2012signatures,rfc8032}; this paper does not prove that
primitive, its library implementation, key custody or a concrete
serialization/refinement theorem. The addition is an explicit standard
security argument and a list of necessary bindings, not cryptographic
novelty or a claim of production security.
